\documentclass[11pt,a4paper]{article}
\usepackage[a4paper,margin=25mm]{geometry}
\usepackage[T1]{fontenc}
\usepackage{amsmath,amssymb,booktabs,tabularx,array,graphicx,xcolor}
\usepackage{hyperref,caption,enumitem,listings,microtype}
\hypersetup{colorlinks=true,linkcolor=black,citecolor=black,urlcolor=blue}
\captionsetup{font=small,labelfont=bf}
\setlength{\parindent}{0pt}
\setlength{\parskip}{4pt}
\setlength{\emergencystretch}{2em}
\setcounter{tocdepth}{1}
\newcolumntype{Y}{>{\raggedright\arraybackslash}X}
\definecolor{draftblue}{RGB}{35,70,115}
\newcommand{\pending}[1]{\textcolor{draftblue}{\textbf{[Insert: #1]}}}
\newcommand{\code}[1]{\texttt{\detokenize{#1}}}
% Replace a frame with \includegraphics, preserving its caption and label.
\newcommand{\figslot}[3]{%
  \begin{figure}[tbp]\centering
  \fbox{\parbox[c][24mm][c]{0.94\linewidth}{\centering\small #3}}
  \caption{#2}\label{#1}\end{figure}}
\lstset{basicstyle=\ttfamily\footnotesize,breaklines=true,columns=fullflexible,
  frame=single,showstringspaces=false}

\begin{document}
\begin{titlepage}
\centering
\vspace*{14mm}
{\Large\bfseries NATIONAL UNIVERSITY OF SINGAPORE\par}
\vspace{4mm}
{\large College of Design and Engineering\par}
{\large Department of Electrical and Computer Engineering\par}
\vspace{15mm}
{\Large\bfseries CEG5206\par}
{\large Algorithms \& Performance Optimization for Computing\par}
\vspace{14mm}
{\LARGE\bfseries FlashAttention Forward\par}
\vspace{4mm}
{\large Materialisation, Causal Fusion and Tensor Core Trade-offs\par}
\vspace{14mm}
{\large Group 1\par}
\vspace{5mm}
\begin{tabular}{ll}
Lyu Sihan & \pending{student ID}\\[3pt]
Zhang Yichun & \pending{student ID}\\[3pt]
Kyaw Tun Khine & \pending{student ID}
\end{tabular}
\vfill
Academic Year 2026/2027, Semester 1\\
13 November 2026
\vspace{8mm}

\small\textbf{Draft status.} Technical content reflects the current A/C implementation.
Part B implementation details and all final experimental findings remain to be supplied.
Figure frames and result slots are intentional; no measured outcomes are asserted.
\end{titlepage}

\begin{abstract}
Conventional scaled dot-product attention writes quadratic score and probability
matrices to GPU global memory. This project studies an alternative decomposition
that computes attention tiles on chip and combines them using numerically stable
online softmax. The implementation comprises materialised dense/causal CUDA
baselines, a dense fused component to be supplied by Part B, and causal fused
CUDA-core and Tensor Core variants. The causal progression explores thread and
warp ownership, shared-memory reuse, two-dimensional register tiling, future-tile
skipping and conditional vector loads. The Tensor Core extension includes both
WMMA and explicit PTX MMA, with its probability-rounding and accumulation contract
stated separately from FP32. A shared FP64 C++ oracle, deterministic input
fingerprints, bounded configuration selection and standardized result tables make
same-machine comparisons reproducible. Evaluation distinguishes latency, explicit
allocation, measured traffic, instruction scheduling and numerical error, while
separating controlled ablations from changes that combine several mechanisms.
Final plots and findings are intentionally reserved for a fresh common-machine
experiment after Part B is integrated.
\end{abstract}
\tableofcontents
\clearpage

\section{Problem, scope and contribution}\label{sec:scope}
For each batch and attention head, the forward operation is
\begin{equation}
 S=QK^{\mathsf T}/\sqrt d+M,\qquad
 P_{ij}=\frac{\exp(S_{ij}-m_i)}{\sum_k\exp(S_{ik}-m_i)},\qquad
 O=PV,\quad m_i=\max_j S_{ij}.
 \label{eq:attention}
\end{equation}
Here $Q,K,V,O$ have shape $[B,H,N,d]$. Dense attention has $M=0$;
causal attention has $M_{ij}=0$ for $j\leq i$ and $-\infty$ otherwise.
These modes have different outputs and different valid interaction counts.
Neither kernel fusion nor the online recurrence approximates the mathematical
operation, although finite precision and reduction order change rounding.

The research question is how restructuring attention around the GPU memory
hierarchy changes performance, and under what shapes the gains justify the
additional on-chip state and control. The study separates three questions:
\begin{enumerate}[nosep,leftmargin=*]
\item How much improvement comes from optimising the materialised pipeline itself?
\item What changes when the score/probability matrices are eliminated, and when
causal structure permits skipping future tiles?
\item How do register/shared-memory constraints and mixed-precision matrix
instructions change the latency--accuracy trade-off?
\end{enumerate}

The required scope is forward FP32 materialised, dense fused and causal fused
attention, with correctness, bounded tuning and reproducible evaluation. Mixed
precision is an extension, not a substitute for a missing dense implementation.
Backward propagation, dropout, arbitrary attention masks, KV-cached decoding,
framework integration and end-to-end training are outside the evaluated scope.
Public names such as FA1/FA2 identify educational variants in this repository;
they do not imply parity with official FlashAttention releases.

Part A supplies an independent FP64 reference and separate CUDA $QK^{\mathsf T}$,
softmax and $PV$ kernels. Part C supplies causal scalar/warp/tiled variants and
fused WMMA/PTX variants. Shared host-side integration connects their tests,
inputs, measurement conventions and analysis without copying GPU kernels between
parts. Part B's dense fused kernel and adapter remain pending in this draft.
The contribution is implementation and evidence-based evaluation, rather than
a new attention algorithm or a claim of production-library performance.

\section{Algorithm and resource model}\label{sec:algorithm}
\subsection{Layout and materialisation}
Tensors use contiguous row-major layout with element offset
$((bH+h)N+i)d+r$. The key tensor is physically $[N,d]$ per head;
$K^{\mathsf T}$ denotes the mathematical transpose. Materialised $S/P$ have
shape $[B,H,N,N]$, with key index contiguous. Strides count elements, not bytes.

For storage width $s$, retaining all six $Q,K,V,S,P,O$ arrays requires
\begin{equation}
 A_{\mathrm{materialised}}=4BHNds+2BHN^2s.
 \label{eq:allocation}
\end{equation}
The quadratic term is auxiliary storage. A write and a subsequent read of both
intermediates amount to $4BHN^2s$ logical bytes, before additional accesses.
This is an access-accounting model, not an unconditional lower bound on DRAM
traffic: cache reuse and the softmax implementation affect physical transfers.
A fused implementation avoids global $S/P$ but still rereads $K/V$ across query
blocks. Its explicit input/output allocation can be linear in $N$ even while
its total arithmetic and data movement remain superlinear.

\subsection{Online softmax and its invariants}
Tiling and online normalisation underpin the FlashAttention algorithm
\cite{flashattention,online}. For one query row, maintain a running maximum $m$,
normaliser $\ell$ and unnormalised output vector $a$. Given a new score tile
$T$, merge it by
\begin{align}
 m'&=\max\bigl(m,\max_{j\in T}s_j\bigr), &
 \alpha&=\exp(m-m'),\nonumber\\
 p_j&=\exp(s_j-m'), &
 \ell'&=\alpha\ell+\sum_{j\in T}p_j,\label{eq:online}\\
 a'&=\alpha a+\sum_{j\in T}p_jv_j, &
 o&=a'/\ell'.\nonumber
\end{align}
This recurrence follows by expressing old and new weights relative to the same
maximum. It preserves the maximum, exponential sum and weighted-value sum over
all processed valid keys. Rescaling only $\ell$ while leaving $a$ unchanged
violates that invariant. Initialise $(m,\ell,a)=(-\infty,0,0)$ and guard the
empty-state scale; an entirely masked subtile must not evaluate
$-\infty-(-\infty)$. Masked weights contribute zero. Some educational variants
normalise the partial output after each update; they must apply the equivalent
rescaling consistently. The unnormalised form allows final division once per row.

\begin{lstlisting}[caption={Conceptual fused forward pass; ownership differs by implementation.},label={lst:online}]
for each independent batch/head and query tile:
    initialise running maximum, denominator and output numerator
    for each retained key/value tile:
        load operands and form score subtiles
        mask out-of-range and future keys
        merge valid scores using the online recurrence
        accumulate probability-times-value contributions
        synchronise before overwriting shared operands
    normalise and store valid output elements
\end{lstlisting}

\subsection{Causal work and tile bounds}
The exact valid interaction counts are
\begin{equation}
 W_{\mathrm{dense}}=BHN^2,\qquad
 W_{\mathrm{causal}}=BH\frac{N(N+1)}2.
 \label{eq:pairs}
\end{equation}
Counting two FLOPs per multiply--add, the two products contribute approximately
$4dW$ useful FLOPs, excluding softmax and control. This numerator is an
algorithmic work model, not the number of instructions executed by padded or
masked kernels.

For query tile $[q_0,q_{\mathrm{end}})$, where
$q_{\mathrm{end}}=\min(N,q_0+B_Q)$, a key tile starting at
$k_0\geq q_{\mathrm{end}}$ is entirely future-facing. Retaining only
$\lceil q_{\mathrm{end}}/B_K\rceil$ key tiles skips such blocks. Remaining
diagonal blocks still need element-wise $k\leq q$ masking. All warps in a CTA
use the same loop bound so that CTA barriers remain collective. This formula
assumes queries and keys share the same position origin; it does not directly
cover offset KV-cache queries.

When $N$ is divisible by $B_Q=B_K=b$, let $T=N/b$. An unskipped tiled schedule
visits $T^2$ block pairs, whereas the skipped schedule visits $T(T+1)/2$.
The asymptotic saving approaches one half of block visits, not necessarily one
half of runtime. Diagonal padding, load imbalance, synchronisation and non-matrix
work remain. Query tiles near the sequence start also perform less work than
those near the end.

\section{Implementation and optimisation decisions}\label{sec:implementation}
\subsection{Part A: materialised comparison baselines}
The naive FP32 pipeline assigns score/output elements to threads and uses a
serial row softmax. The tiled matrix kernels cooperatively stage operands in
shared memory and accumulate output elements in registers. The $QK^{\mathsf T}$
kernel pads the shared K tile by one column; this layout choice must be evaluated
rather than assumed to eliminate every conflict. Parallel softmax assigns a
block to a row, combines per-thread maxima and sums through power-of-two shared
reductions, and writes normalised probabilities. Separate variants expose tiled
matrix multiplication and parallel softmax independently before combining them.

\begin{table}[tbp]\centering\small
\begin{tabularx}{\linewidth}{@{}lYY@{}}\toprule
Variant & Matrix products & Softmax and purpose\\\midrule
\code{naive} & Direct FP32 kernels & Serial row baseline\\
\code{tiled} & Shared-memory FP32 tiles & Serial; isolates product tiling\\
\code{parallel} & Direct FP32 kernels & Block reduction; isolates softmax mapping\\
\code{optimized} & Shared-memory FP32 tiles & Block reduction; stronger fusion comparator\\
\code{tensorcore} & Separate FP16 WMMA GEMMs & FP16 intermediates, reduction arithmetic and output\\
\bottomrule\end{tabularx}
\caption{Part A variants. The Tensor Core path changes numerical precision as well as matrix hardware.}\label{tab:parta}
\end{table}

Both dense and causal modes retain the materialised buffers. Optimised causal
$QK^{\mathsf T}$ can avoid multiply work for fully masked tiles while still
writing masked scores; therefore A-causal is not necessarily a full dense
computation with a mask added afterwards. It remains an appropriate same-mode
reference for C. Stage measurements use separate passes: their medians must not
be summed to reconstruct the independently timed full forward pass.

\subsection{Part B: dense fused component}
The intended dense kernel assigns query tiles to CTAs, visits all required K/V
tiles and applies Equation~\ref{eq:online} without causal masking. It must avoid
global $S/P$, validate against the shared dense FP64 oracle, and freeze a bounded
configuration search before the reporting grid. This is the integration contract,
not a description of an already supplied implementation.

\pending{B author: add the actual kernel progression, CTA/warp ownership,
on-chip layouts, bounded candidates and any implementation-specific limitations.
Keep this to one focused subsection and update its status after validation.}
The final report must compare B-dense against A-dense; comparison with C-causal
instead addresses a change in workload. B publishes the same seven tables and
runner interface described in Section~\ref{sec:reproducibility}.

\subsection{Part C: CUDA-core progression}
\begin{table}[tbp]\centering\small
\begin{tabularx}{\linewidth}{@{}lYY@{}}\toprule
ID & Organisation & Principal design change\\\midrule
\code{fa1} & One thread per query row & Educational online-softmax baseline; no shared staging\\
\code{scalar} & Scalar query-row work & FA2-style recurrence variant; not an official FA2 implementation\\
\code{warp} & Warp per query & Feature-parallel reductions and shared K/V reuse\\
\code{tiled} & Multiple query rows per warp & Shared Q/K/V and 2D register score tiles\\
\code{skip} & Tiled, causal-bounded traversal & Avoid entirely future K/V tiles\\
\code{vector} & Same optimised family & Conditional \code{float4} staging with scalar fallback\\
\bottomrule\end{tabularx}
\caption{Maintained causal CUDA-core variants. Progression includes combined changes and is not automatically a sequence of isolated ablations.}\label{tab:partc}
\end{table}

In the tiled family, grid x selects a batch/head pair and grid y a query tile.
A CTA has $32W$ threads. Warp $w$ owns up to $R=\lceil B_Q/W\rceil$ query rows,
with lane $\lambda$ owning feature indices $32g+\lambda$. Bounds predicates cover
partial query, key and feature tiles. Q/K/V shared staging has logical payload
\begin{equation}
 A_{\mathrm{shared,core}}=4d(B_Q+2B_K)\ \text{bytes}.
 \label{eq:shared}
\end{equation}
Barriers separate completed loads, tile consumption and overwriting K/V storage.
The public wrappers enforce a contiguous layout even where a low-level kernel
has more general stride parameters; arbitrary-stride coverage is not claimed.

A thread's $T_Q\times T_K$ partial score tile reuses query/key register operands
across several dot products. Lanes reduce partial dot products through warp
shuffles. Running row maxima and denominators are replicated across lanes,
while output features are distributed. Larger register tiles can reduce repeated
loads but increase live state, dependency chains and local-memory traffic.
Source arrays called register tiles are not proof that the compiler keeps every
value in registers.

For vector staging, a lane loads four consecutive features when source alignment,
contiguity and bounds permit a 16-byte transaction; otherwise it uses guarded
scalar accesses. This loading ownership differs from arithmetic ownership.
Output stores remain scalar and coalesced across lanes. Wider accesses can lower
instruction count without lowering DRAM bytes, because adjacent scalar accesses
may already coalesce. Consequently vectorisation is a measured option rather
than an assumed final winner.

\subsection{Part C: WMMA and explicit PTX MMA}
Both Tensor Core variants fuse score formation, online softmax and value
aggregation into the forward kernel. They take FP16 Q/K/V and retain FP32
accumulation, softmax state and output; probability operands are rounded to FP16
before $PV$. Their oracle uses the same FP16-rounded inputs, while tolerance
also accommodates internal arithmetic and probability rounding.

The WMMA implementation performs $16\times16\times16$ warp-level products,
stores score fragments to shared memory for masking/softmax, and stages probability
and product tiles for the second matrix operation. This is already a fused Tensor
Core implementation: direct PTX is not a requirement for functional fusion.

The PTX variant issues
\code{mma.sync.aligned.m16n8k16.row.col.f32.f16.f16.f32} with explicitly packed
half operands. A warp owns 16 query rows. Two $16\times8$ score products cover
16 keys, after which four-lane groups reduce the corresponding row scores.
FP32 output fragments are rescaled by the online correction and reused for
$PV$; probabilities pass through shared half storage. Accumulator mapping permits
direct FP32 output stores without an output shared-memory tile. This exposes
fragment ownership and avoids some WMMA score/output staging, while adding
manual packing and mapping complexity. The instruction contract is defined by
NVIDIA's PTX ISA \cite{ptx}.

The current PTX path requires $d\geq16$ divisible by 16; WMMA pads internally.
Supported wrappers cover a bounded instantiation set, not arbitrary dimensions.
The maintained C kernels use neither CUTLASS/CuTe nor \code{ldmatrix}, shared-memory
swizzling, asynchronous-copy pipelines or double buffering. These are possible
future techniques, not implemented optimisations or explanations for existing gains.

\begin{table}[tbp]\centering\small
\begin{tabularx}{\linewidth}{@{}lYYY@{}}\toprule
Path & Inputs & Accumulation/state & Stored intermediates/output\\\midrule
A FP32 & FP32 & FP32 & FP32 S/P/O\\
B FP32 contract & FP32 & FP32 & FP32 O; no global S/P\\
C CUDA core & FP32 & FP32 & FP32 O; no global S/P\\
A Tensor Core & FP16 & FP16 GEMM/softmax arithmetic & FP16 S/P/O\\
C WMMA/PTX & FP16 & FP32 & FP16 probability operands; FP32 O\\
\bottomrule\end{tabularx}
\caption{Numerical contracts. Equal FP16 inputs do not make A/C Tensor Core paths precision-equivalent. B remains a required interface until its code is supplied.}\label{tab:precision}
\end{table}

\section{Correctness and integration}\label{sec:correctness}
\subsection{Shared independent oracle and input identity}
Part A's reusable C++ \code{attention_reference} library provides full materialised
FP64 S/P/O and an output-only FP64 reference. The latter evaluates either all
queries or the unique set $\{0,\lfloor N/2\rfloor,N-1\}$ in each head, without
allocating full score/probability matrices. Agreement tests compare both helpers
with the original reference in dense and causal modes, including tails and
structured patterns. CPU causal-isolation tests perturb future K/V and check
that preceding outputs are unchanged; they establish a reference invariant and
do not replace GPU output tests.

All current standardized A/C runners call the same SplitMix64 generator,
seed 5206, filling Q then K then V. FNV-1a-64 fingerprints cover canonical
little-endian FP32 bytes in that order. FP16 paths also fingerprint rounded
values represented as FP32. Shared pattern transformations avoid independent
implementations silently generating different tensors. Equal seed labels alone
are insufficient. Historical records preserve their original generator/oracle
provenance and are not relabelled as fresh shared runs.

\subsection{Checks, tolerances and limits}
For each checked element the acceptance criterion is
\begin{equation}
 |x_i-r_i|\leq a_{\mathrm{tol}}+r_{\mathrm{tol}}|r_i|.
 \label{eq:tolerance}
\end{equation}
FP32 uses $a_{\mathrm{tol}}=r_{\mathrm{tol}}=10^{-4}$; current FP16-input experiments
use $a_{\mathrm{tol}}=0.002$, $r_{\mathrm{tol}}=0.02$. Maximum absolute and relative
errors are recorded, with the latter using denominator $\max(|r_i|,10^{-12})$.
Near-zero references can produce large relative maxima despite passing the
combined criterion. Tolerances are fixed before results, not widened after failure.
Expected masked $S=-\infty$ entries are treated separately from nonfinite output.

\begin{table}[tbp]\centering\small
\begin{tabularx}{\linewidth}{@{}lY@{}}\toprule
Test family & Failure mechanism addressed\\\midrule
N=1, tails and multiple heads & Boundary predicates, partial tiles and batch/head offsets\\
Uniform attention and constant V & Normalisation and weighted-sum invariants\\
Rising/stress logits & Stable exponentials and rescaling when later tiles change maxima\\
Small FP64 comparisons & Full A S/P/O and fused O errors, including diagonal masking\\
Candidate configuration checks & Invalid warp/tile mappings before selection\\
Per-benchmark output validation & Exact config/input linkage; all-output finiteness\\
Compute Sanitizer & Representative memcheck, racecheck, synccheck and initcheck\\
Schema/analysis regression & Missing cases, duplicate samples, invalid statistics and provenance\\
\bottomrule\end{tabularx}
\caption{Validation design. Representative sanitizers and sampled large outputs do not establish exhaustive coverage.}\label{tab:validation}
\end{table}

Small correctness suites compare full outputs. Part A also checks all O for its
small performance shapes; large A cases and C benchmark cases use sampled
queries plus finiteness checks over every O element. Saved scope, query indices
and checked-element counts distinguish these guarantees. Unsupported PTX
shapes are explicit skips, not passing executions. Every accepted timed case
must link to passing O validation outside the timing interval.

The CPU oracle is authoritative. The proposal additionally calls for fused
versus materialised GPU output comparisons; these are a separate final acceptance
check on identical dense or causal inputs. Agreement of timing CSVs cannot prove
output equality, and the shared-oracle smoke audit alone is not evidence of
that additional direct GPU-to-GPU check.

\section{Experimental methodology}\label{sec:method}
\subsection{Platform, grid and selection}
The final comparison uses one physical GPU and one common run. Source/binary
hashes, GPU UUID, compiler/driver versions, architecture, commands and dependencies
are retained. Earlier A and C measurements on different devices are separate
experiments and are excluded from cross-part speedup denominators.

\begin{table}[tbp]\centering\small
\begin{tabularx}{\linewidth}{@{}lY@{}}\toprule
Setting & Full experiment contract\\\midrule
Input shape & $B=4$, $H=8$; $N\in\{256,512,1024,2048,4096\}$; $d\in\{32,64,128\}$\\
Modes & A dense/causal; B dense; C causal\\
Core/extension & FP32 required; FP16 paths separately labelled\\
Tuning & $N=256,1024$; selection frozen per head dimension (and A mode)\\
Held-out lengths & $N=512,2048,4096$ are excluded from selection\\
Timing & Five warmups; 20 CUDA-event forward samples; median, Q1/Q3 and IQR\\
Profiling & Separate frozen $d=64$, $N=256/4096$ captures, three repetitions\\
Final platform & \pending{GPU/VRAM, architecture, driver, CUDA/nvcc, host, Python/package versions}\\
\bottomrule\end{tabularx}
\caption{Common measurement plan. The small smoke grid checks integration and is not report benchmark evidence.}\label{tab:grid}
\end{table}

A searches FP32 $(\mathrm{tile},\mathrm{softmax\ threads})\in
\{(8,256),(16,256),(32,256),(16,128)\}$ and FP16 $(16,128)/(16,256)$.
C searches five $(B_Q,B_K,W,T_Q,T_K)$ configurations:
$(32,32,4,3,5)$, $(32,32,4,2,4)$, $(16,32,4,2,4)$,
$(32,16,4,2,4)$ and $(32,32,8,2,4)$ for skip/vector;
WMMA/PTX use $(B_Q,B_K,W)\in\{(32,32,2),(16,32,1),(32,16,2),(16,16,1)\}$.
Selection minimises geometric mean median latency across tuning lengths, with
deterministic tie-breaking. C first freezes each variant, then selects a family
winner among skip/vector or WMMA/PTX. Earlier scalar/warp variants remain fixed
progression references; family selection is not an exhaustive search over all
possible implementations. B must state its own bounded candidates.

Tuning lengths intentionally overlap the report grid. Results at those lengths
measure tuned performance; only the remaining lengths test generalisation to
held-out sequence lengths. No reporting-grid minima are selected retrospectively.

\subsection{Timing and metrics}
Allocation, generation, rounding, transfers, CPU reference, resource queries and
CSV writing are outside the CUDA-event interval. A full forward includes all
three kernels; C includes both products and online softmax in one kernel.
Profiler-instrumented timings are excluded from accepted samples. The runners
execute workloads sequentially on the selected device. They do not randomise
order or lock benchmark clocks, so thermal/cache history remains a limitation.

For forward median $t$ milliseconds, sequence throughput is $1000B/t$ and token
throughput is $1000BN/t$. Same-operation speedup is $t_{\mathrm{baseline}}/
 t_{\mathrm{candidate}}$, with shape, mode and numerical contract stated. These
are attention-kernel metrics, not complete model inference/training throughput.
Twenty samples describe within-case spread; they do not provide independent-session
confidence intervals. A single geometric mean across shapes, if used, must
state the included cases and must not replace the per-shape plots.

Explicit memory is the live allocation ledger, including Q/K/V/O and materialised
S/P where applicable. C uses 16 bytes per tensor element for FP32 Q/K/V/O,
or 10 for FP16 Q/K/V and FP32 O. These exclude context/runtime allocations.
Observed free-memory deltas at a few checkpoints are supplementary, not true
process high-water marks. Profiler DRAM/local/shared transactions are traffic,
not capacity. Static resource occupancy estimates are likewise distinct from
measured active-warp occupancy.

\subsection{Controlled comparisons and attribution}
The principal same-mode comparisons are optimized A-dense versus B-dense and
optimized A-causal versus C-causal. C CUDA-core versus C Tensor Core describes
a mixed-precision implementation trade-off. A-Tensor-Core versus C-Tensor-Core
has equal rounded inputs but different accumulation/intermediate/output precision
(Table~\ref{tab:precision}); it cannot isolate fusion alone.

A's separate tiled-only and parallel-only variants help attribute baseline gains.
For C, the matched skip/vector study uses identical launch parameters and
arithmetic/output mapping, changing scalar versus conditional vector loading.
By contrast, tiled-to-skip and WMMA-to-PTX may combine several changes. They
are progression comparisons, not clean single-feature causal experiments.
Default-versus-frozen comparisons show configuration sensitivity. Dense-B versus
causal-C changes the valid interaction set and possibly kernel organisation;
its latency ratio is not a speedup for identical computation.

\section{Results to insert and analysis framework}\label{sec:results}
No empirical values, winners or bottleneck findings are populated in this draft.
All slots below refer to the final accepted common-machine bundle; historical
figures may be discussed only as separately labelled supplementary experiments.
Captions define the required evidence so that inserting a plot does not require
repeating methodology in each subsection.

\subsection{Validation and implementation progression}
\pending{Insert final run ID/commit, tested and unsupported case counts by part,
full/sample scope, maximum absolute/relative errors by precision, and sanitizer
outcomes. Confirm B completion and the supplementary materialised-GPU checks.}
A passing aggregate audit is necessary but does not turn sampled checks into
full large-output validation. Inspect errors as $N/d$ grow and identify any
precision-specific restriction before interpreting speedups.

\figslot{fig:cuda-progression}{CUDA-core progression across sequence length and head dimension. Separate A dense/causal and C causal panels; report medians with IQR and identify fixed-default versus frozen settings. B panels are added after its kernel progression is supplied.}{Insert FP32 progression panels from the final per-part plots.}
\figslot{fig:tensor-progression}{C WMMA/PTX progression and frozen CUDA-core versus Tensor Core comparison. Label FP16 inputs, FP32 state/output and probability rounding; include compatible A context with its different precision contract.}{Insert Tensor Core progression and frozen-family comparison panels.}

Progression plots answer where a design change helps and where it loses. Discuss
whether effects persist across head dimensions, whether a winner is stable
relative to IQR, and whether an apparent gain survives held-out lengths. A
slower vector variant is a valid result: instruction reduction, transaction size,
register pressure and launch configuration can act in opposite directions.

\subsection{Materialised versus fused and memory scaling}
\figslot{fig:fusion}{Same-input materialised/fused comparisons: A-dense/B-dense and A-causal/C-causal. Show median forward latency, IQR and explicitly identified speedup denominators; do not mix historical GPUs.}{Insert the final combined latency/speedup panels.}
\figslot{fig:memory}{Explicit live allocation and throughput scaling for matched shapes. Keep analytical S/P storage, measured allocation ledgers and profiler DRAM bytes visibly distinct.}{Insert allocation/throughput panels and, where available, measured traffic.}

The allocation model predicts a quadratic-versus-linear capacity distinction,
not a universal latency crossover. For short sequences, launch overhead and
extra fused-kernel bookkeeping may dominate. For longer sequences, reuse and
avoided intermediates may matter more, but each query CTA still streams K/V.
Use traffic counters and resource evidence to determine whether reduced capacity
also corresponds to reduced off-chip traffic. Larger $d$ increases output state
and on-chip storage, potentially shifting the crossover. Attribute the observed
pattern only after inspecting the measurements.

\subsection{Tuning and controlled loading study}
\figslot{fig:tuning}{Bounded tuning and default/frozen sensitivity by head dimension. Record selected configurations before the final grid; distinguish tuning lengths from held-out lengths.}{Insert tuning heatmap and default-versus-frozen panels.}
\figslot{fig:load-ablation}{Matched scalar/vector loading with identical launch configuration and compute mapping. Report skip/vector latency ratios and IQR, including slowdowns.}{Insert the matched-load ablation from the accepted C study.}

A larger tile raises reuse but also shared-memory and register requirements.
Use Equation~\ref{eq:shared}, compiler/resource queries and measured latency to
explain a configuration choice. Do not equate lower local bytes per thread with
fewer executed local-memory transactions. The matched experiment can support a
claim about loading at its tested launch; it does not establish vectorisation
as universally better or explain every difference in the broader progression.

\subsection{Dense versus causal work}
\figslot{fig:causal-work}{Dense B versus causal C on identical base inputs: forward latency, analytical valid-pair counts and compatible measured traffic. These are different workloads, not an output-equivalence comparison.}{Insert dense/causal workload panels after B is available.}

Compare interaction reduction in Equation~\ref{eq:pairs} with actual skipped
block schedules, accounting for diagonal over-computation. A smaller valid-pair
count can coexist with imperfect occupancy, uneven CTA duration and unchanged
launch/storage overhead. An analytical tile count is not measured executed
instruction count. If a near-twofold latency reduction occurs, establish whether
it follows work elimination or also reflects different tiling/precision.

\subsection{Profiler evidence and remaining bottlenecks}\label{sec:analysis}
\figslot{fig:profile}{Representative frozen-kernel evidence at $d=64$, $N=256/4096$: achieved occupancy, eligible warps, sampled warp states and local/shared/DRAM traffic. Keep counter availability and capture repetitions explicit.}{Insert a compact profiler panel; retain full metrics/source tables in the repository.}

A useful analysis follows the chain from source mechanism to emitted instructions,
measured counters and finally unprofiled latency. The roofline relation
$P\lesssim\min(P_{\mathrm{peak}},I\beta)$ provides context, but
$I=\mathrm{useful\ FLOPs}/\mathrm{measured\ DRAM\ bytes}$ alone cannot diagnose
on-chip bandwidth, dependencies, synchronisation or special-function overhead.
It is not proof that an attention kernel is compute-bound whenever DRAM traffic
falls.

\begin{table}[tbp]\centering\small
\begin{tabularx}{\linewidth}{@{}lYY@{}}\toprule
Evidence & Interpretation to test & Required qualification\\\midrule
High local traffic & Register allocation/indexed arrays may cause local accesses & Inspect SASS/source; allocation bytes alone are insufficient\\
Low eligible-warps rate & Dependency, memory or synchronisation limits issue readiness & Occupancy alone does not measure useful issue throughput\\
MIO/short-scoreboard samples & Investigate shared/local operations and instruction dependencies & Sampled states are not runtime percentages; MIO is not a bank-conflict count\\
Excess shared wavefronts & Investigate lane/address layout at the source instruction & Separate actual/ideal wavefronts from aggregate arbitration counters\\
Low Tensor Core activity & MMA may be limited by packing, softmax, staging or dependencies & Instruction presence does not establish end-to-end utilisation\\
Predication/branch metrics & Investigate tails, diagonal masks and lane participation & Predicated-off lanes are not automatically divergent branches\\
\bottomrule\end{tabularx}
\caption{Evidence hierarchy for interpreting the final profiles, rather than pre-written bottleneck claims.}\label{tab:profiling}
\end{table}

NVIDIA documents replay, metric units, scheduler states and counter semantics
\cite{nsight}. Normalise sampled warp states only with an explicit denominator;
zero samples or unavailable counters must not be rendered as zero stalls.
Source-correlated excessive shared wavefronts provide stronger layout evidence
than an aggregate bank-conflict counter alone. If that evidence is absent, do
not claim that swizzling would fix the kernel. Padding already present in A is
also different from an implemented C swizzle.

\pending{Replace this slot with measured explanations: one supported improvement,
one slowdown/limited gain, the principal remaining bottlenecks, and evidence
paths/figure references. State what counters cannot establish.}

\section{Limitations, next steps and conclusion}\label{sec:limits}
The final evaluation is bounded by one selected GPU, a fixed shape grid and
synthetic inputs. Results need not generalise to different architectures,
longer contexts, adversarial distributions or application-level workloads.
Fixed ordering, unmanaged clocks and within-case sampling limit uncertainty
claims. Large-output sampling can miss errors outside the checked rows, and
sanitizers are representative rather than exhaustive. Mixed precision introduces
both rounding and implementation differences, while compiler-managed register
placement complicates source-only reasoning.

The most relevant next experiments follow the actual bottleneck evidence.
Swizzling may address demonstrated shared-layout conflicts; asynchronous copies
and double buffering may overlap transfers with computation, but need correct
producer/consumer scheduling and may increase shared-memory pressure. Reducing
fragment staging or indexed local arrays may be preferable if those dominate.
These are hypotheses, not promised gains, and none is required merely to name
another optimisation. The current goal is completion and evaluation of the
approved dense/causal scope, not an unbounded rewrite of Part C.

The implemented design exposes how materialisation, thread ownership, on-chip
reuse, causal scheduling and matrix instruction interfaces change the execution
of exact attention. Shared validation and provenance make the comparisons
reviewable. \pending{Insert a short final conclusion after experiments: the
conditions where fusion/causal skipping help, the measured precision trade-off,
and the strongest supported limitation. Do not claim a winner before the run.}

\section{Reproduction, contributions and attribution}\label{sec:reproducibility}
The root and standalone runners build through CMake, validate before accepting
timings, freeze configurations, collect evidence and generate tables/plots.
Seven public CSVs share a versioned schema: correctness, samples, summary,
resources, profiles, memory and work. Joins preserve implementation/configuration,
shape, mode, precision, fingerprints and validation scope. Root aggregation
checks the selected grid, physical GPU identity, sample statistics and evidence
links. Native records and supplementary insights remain in each child bundle;
empty fields distinguish unavailable measurements from measured zeros.

\begin{lstlisting}[caption={From the repository root after dependency setup; use fresh output directories.},label={lst:reproduce}]
# Available A/C smoke integration check; not final report data.
bash run.sh --parts a,c --grid smoke --output results/ac_smoke
# Final common-machine experiment, after Part B is supplied.
bash run.sh --output results/final_submission
# Standalone executions use the same public format.
bash part_a/run.sh --output results/a_only
bash part_b/run.sh --output results/b_only  # requires completed B
bash part_c/run.sh --output results/c_only
# CPU-only native re-extraction, audits and figure regeneration.
bash run.sh --analyze-only results/final_submission
\end{lstlisting}

\code{run_all.sh} is an alias of the root runner. Defaults request all parts;
missing B causes preflight failure rather than silently producing a complete-run
claim. Smoke and no-profile runs retain their partial-evidence labels. The
submission's \code{README.txt} contains installation and examiner run instructions;
\code{README.md} specifies B's interface and exact output paths. For a complete
run, aggregate CSVs are at \code{results/<run_id>/}, individual bundles at
\code{parts/part_a/}, \code{parts/part_b/}, \code{parts/part_c/} beneath that root,
and report figures at \code{plots/<part>/} and \code{plots/combined/}.

Commit compact numerical evidence, raw metric CSVs, processing scripts, manifests,
configuration selections and report figures needed to reproduce cited findings.
Large binary Nsight reports may be distributed separately with an accessible link,
checksums and restore paths; CSV-only reproduction does not provide every GUI or
SASS investigation. Avoid distributing duplicate native/profile archives merely
because both native and standardized exports retained a copy.

\begin{table}[tbp]\centering\small
\begin{tabularx}{\linewidth}{@{}lY@{}}\toprule
Member & Responsibility and contribution to confirm before submission\\\midrule
Lyu Sihan & Part A FP64 oracle and materialised GPU baselines; review final baseline evidence.\\
Zhang Yichun & Part B dense fused implementation, testing and bounded tuning; pending in this draft.\\
Kyaw Tun Khine & Part C causal CUDA-core/Tensor Core kernels, evaluation and shared A/C integration.\\
Team & Review reported findings, actual contributions, video, citations and submission artifact.\\
\bottomrule\end{tabularx}
\caption{Assigned responsibilities; update to actual completed contributions.}\label{tab:contributions}
\end{table}

FlashAttention and FlashAttention-2 supply the IO-aware and work-partitioning
context \cite{flashattention,flashattention2}; the implementation is not the
official library. The C educational baseline explicitly cites Kim's minimal CUDA
implementation in its source \cite{minimal}. Li's tutorial/repository was consulted
as an external optimisation comparison \cite{li}; its benchmark percentages and
optimisations are not our experimental results. \pending{Confirm exact adapted
code/lines, licenses and team modifications; retain source-level attribution.}

ChatGPT/Codex assisted code development, host-side integration, testing/analysis
scripts and this report draft. This assistance is not independent experimental
evidence. Team members must verify the generated explanations, understand the
submitted code, review references and approve the reported contribution statement.
\pending{Confirm the final AI-use description with all members.}

\textbf{Submission links.} Repository: \pending{clickable public repository URL};
video: \pending{accessible 10--15 minute video URL}; final revision:
\pending{commit hash or release tag}. Supply any separately hosted profiler archive
with its checksummed manifest and accessible download link.

\begin{thebibliography}{9}
\bibitem{flashattention}
T.~Dao, D.~Y.~Fu, S.~Ermon, A.~Rudra and C.~R\'e,
``FlashAttention: Fast and Memory-Efficient Exact Attention with IO-Awareness,''
\emph{NeurIPS}, 2022. \url{https://arxiv.org/abs/2205.14135}.
\bibitem{flashattention2}
T.~Dao, ``FlashAttention-2: Faster Attention with Better Parallelism and Work
Partitioning,'' 2023. \url{https://arxiv.org/abs/2307.08691}.
\bibitem{online}
M.~Milakov and N.~Gimelshein, ``Online normalizer calculation for softmax,''
2018. \url{https://arxiv.org/abs/1805.02867}.
\bibitem{ptx}
NVIDIA, \emph{Parallel Thread Execution ISA}, matrix instructions and fragment
layouts. \url{https://docs.nvidia.com/cuda/parallel-thread-execution/}.
Record the toolkit/ISA version used in the final run manifest.
\bibitem{nsight}
NVIDIA, \emph{Nsight Compute Profiling Guide}.
\url{https://docs.nvidia.com/nsight-compute/ProfilingGuide/index.html}.
Record the actual capture-tool version with the data.
\bibitem{minimal}
P.~Kim, \emph{flash-attention-minimal}, educational CUDA implementation.
\url{https://github.com/tspeterkim/flash-attention-minimal}.
\bibitem{li}
S.~Li, \emph{Flash Attention from Scratch}, tutorial and companion source.
\url{https://lubits.ch/flash/};
\url{https://github.com/sonnyli/flash_attention_from_scratch}.
\end{thebibliography}
\end{document}
